% Exercise 15.6
We assume the mass of the Ping-Pong ball to be $m_p = 2.7\,\text{g} = 0.0027\,\text{kg}$ 
and its diameter to be $d_p = 38\,\text{mm} = 0.038\,\text{m}$, while those of the bowling
ball are $m_b = 5.4\,\text{kg}$ and $d_b = 21.59\,\text{cm} \approx 0.216\,\text{m}$, respectively.
Thus the cross-sectional of each (approximately) spherical object is given by
\[
    A_p = \pi r_p^2\,\text{m}^2 = \pi \cdot {0.019}^2\,\text{m}^2,
\]
and 
\[
    A_b = \pi r_b^2\,\text{m}^2 = \pi \cdot {0.108}^2\,\text{m}^2.
\]

If we use Equation 14.9 for air resistance with a drag coefficient of $C = 1/2$ and air density
$\rho = 1.225\,\text{kg}/\text{m}^3$, then we have the two equations
\begin{equation*}
    \begin{split}
        F_{\text{air}_p} &= -\frac{1}{2}C\rho A_p \, \lvert v_p \rvert \, v_p \\
                         &= -\frac{1.225\pi \cdot {0.019}^2}{4} \, \lvert v_p \rvert \, v_p,
    \end{split}
\end{equation*}
and
\begin{equation*}
    \begin{split}
        F_{\text{air}_b} &= -\frac{1}{2}C\rho A_b \, \lvert v_b \rvert \, v_b \\
                         &= -\frac{1.225\pi \cdot {0.108}^2}{4} \, \lvert v_b \rvert \, v_b.
    \end{split}
\end{equation*}

In Figure~\ref{fig15_6_i}, we have a graph of velocity as a function of time of the Ping-Pong 
ball and bowling ball dropped from a height of $100\,\text{m}$, and in Figure~\ref{fig15_6_ii}
we have the same from a height of $500\,\text{m}$.

\begin{figure}[!ht]
    \includegraphics[width=\textwidth, center]{../plots/15_6_i.pdf}
    \caption{Velocity as a function of time after drop from $100\,\text{m}$ \label{fig15_6_i}} 
\end{figure}

\begin{figure}[!ht]
    \includegraphics[width=\textwidth, center]{../plots/15_6_ii.pdf}
    \caption{Velocity as a function of time after drop from $500\,\text{m}$ \label{fig15_6_ii}} 
\end{figure}
